\section{An exact local boundary-distribution criterion for the full source}
\label{sec:local-boundary-distribution-criterion}

\paragraph{Status.}
The complete author and two independent whole-proof readings
have passed.  Their exact scopes are recorded separately.
No arithmetic validity of the equivalent conditions is asserted.

Use the original source settings of (MA-MO.700): a nontrivial
primitive finite-order Hecke character $\eta$ with its dual,
structural $S$, and actual Fourier phases; or the rational principal
sector with $S=\{\infty\}$ and its two actual polar sections.
Keep the original real smooth step $\phi$, equal to one far left
and zero far right, with compactly supported derivative; put
$q=1-\phi$.
Write
\[
\begin{gathered}
 H=H_{\Sigma,a,s},\qquad
 w(x)=(e^{2ax}+e^{-2ax})(1+x^2)^r,\qquad r=s/2,\\
 0<a<1/2,\qquad s\ge0,\qquad
 0<\epsilon_0<\tfrac12(1/2-a),\\
 A_\omega(z)=E_{\omega,S}(1/2+z)L_K(1/2+z,\omega)
       \quad(\omega=\eta,\bar\eta),\\
 A(z)=\zeta(1/2+z)\quad\hbox{in the rational case},\qquad
 \mathcal W=H\oplus\mathbb C^2 .
\end{gathered}
\tag{MA-MO.1500}
\]
The norm of $H$ is exactly $L^2(w\,dx)$.  The two polar
coordinates of $\mathcal W$ are the actual section coordinates
of MA-MO-016; no bare polar function is put in $H$.

Throughout this theorem assume the closed-strip premise
\[
 |\Re(\rho-1/2)|\le a
 \quad\hbox{for every completed nontrivial zero, in both
 dual sectors when applicable}.
\tag{MA-MO.1501}
\]
No individual boundary-retention hypothesis is added.
The finite Euler factors have no zeros on the positive
central strip used below.  The rational pole at $z=1/2$
is outside $a\le\Re z\le a+\epsilon_0$.

Choose a smooth uniform partition of unity on the frequency
line: for some $\ell>0$, $t_n=n\ell$ and
$\chi_n(u)=\chi(u-t_n)$, where $\chi\in C_c^\infty$,
the supports have uniformly bounded overlap, and
$\sum_{n\in\mathbb Z}\chi_n=1$.  The choice can, for example,
have $\operatorname{supp}\chi\subset[-2\ell,2\ell]$.
Define the following condition, denoted $\mathrm{LB}_{a,s}$:
for every indicated $\omega$ there are
$d_{\omega,n,0}\in H^{-r}(\mathbb R)$ and constants $C,D$
such that
\[
\begin{aligned}
 d_{\omega,n,\epsilon}(u)
   &:=\frac{\chi_n(u)}{A_\omega(a+\epsilon+iu)}
       \longrightarrow d_{\omega,n,0}(u)
                &&\text{strongly in }H^{-r}\text{ for each }n,\\
 \|d_{\omega,n,\epsilon}\|_{H^{-r}}
   &\le C(2+|n|)^D
                &&(n\in\mathbb Z,\ 0<\epsilon\le\epsilon_0).
\end{aligned}
\tag{MA-MO.1502}
\]
The rational condition has the one self-dual multiplier $A$.
The Sobolev norm uses the unitary Fourier transform in $u$.
There is no uniform convergence rate in $n$ in (MA-MO.1502).
Existence for one such partition is equivalent to existence
for every fixed smooth uniform partition, as proved below.

\paragraph{MA-MO-036 (the actual universal bounded family has an exact boundary criterion).}
Under (MA-MO.1501), the following three conditions are equivalent.
In the nontrivial sector they are:
\[
\begin{array}{ll}
{\rm (B)}&
 \mathscr D_{\eta,g}^{\,\mathrm{raw}}
   =C_g(P_+-P_-):\mathcal E\to\mathcal E
 \text{ has a bounded extension on }H\\
&\hfill
 \text{for every }g\in C_c^\infty([-2,-1]\cup[1,2]);\\
{\rm (L)}&\mathrm{LB}_{a,s}\text{ holds for both }\eta,\bar\eta;\\
{\rm (S)}&
 \mathscr D_{\eta,g}^{\,\mathrm{raw}}\text{ has a unique }
 \mathcal S_1(H)\text{ extension for every }g\in C_c^\infty,\\
&\hfill\text{continuously in trace norm on each fixed support}.
\end{array}
\tag{MA-MO.1503}
\]
In the rational principal sector replace (B) by the weaker
negative-column assertion
\[
\begin{gathered}
 \mathscr D_g^{\,\mathrm{raw}}i_-:\mathcal E\to\mathcal W
 \text{ has a bounded extension }H\to\mathcal W\\
       \text{for every }g\in C_c^\infty([1,2]),
\end{gathered}
\tag{MA-MO.1504}
\]
and replace (S) by the full original matrix extension
$\mathscr D_g\in\mathcal S_1(\mathcal W)$, with the same
continuity.  In both settings (B) can be required only for
real tests, without changing the equivalence.

Whenever these conditions hold, every boundary zero of
multiplicity $m$ satisfies
\[
                         s>2m-1 .
\tag{MA-MO.1505}
\]
Thus the actual full-retention, relative-trace and determinant
conclusions of MA-MO-029 apply, with their unchanged primitive
Euler factors, polar correction and $g(0)=0$ restriction for
the stated prime formulas.  This is an analytic equivalence
under (MA-MO.1501), not a proof of either premise for an
arithmetic divisor.

\begin{proof}
\emph{1. Universal boundedness gives a finite test seminorm.}
Use $\mathcal D_+=C_c^\infty([1,2])$ with its closed
Fr\'echet topology.  In the nontrivial case let
$T_g=\mathscr D_{\eta,g}$; in the rational case let
$T_g=\mathscr D_g i_-$, with output space $\mathcal W$.
For a fixed compact smooth input $u$, $g\mapsto T_g^0u$
is continuous in the indicated output norm.  In the first
case this follows from $T_g^0u=C_gD^0u$ and the weighted
translation bound, since $D^0u\in\mathcal E$.  In the
second case it follows from the actual strongly continuous,
locally bounded representation on $\mathcal W$,
(MA-MO.259)--(MA-MO.261), including its cocycle and both
polar characters.  No bounded full matrix is assumed here.

The closed sets
$\{g:\|T_g^0u\|\le j\|u\|_H\text{ for all compact smooth }u\}$
cover $\mathcal D_+$ under (B).  Baire's theorem, subtraction,
linearity and scaling therefore give
\[
              \|T_g\|\le C\|g\|_{C^N}
                         \qquad(g\in\mathcal D_+)
\tag{MA-MO.1506}
\]
for some finite $N$.  The original compact core is dense,
so each extension is unique.  If (B) is assumed only for
real tests, apply Baire in the real Fr\'echet space first
and then write a complex test as its real and imaginary
parts.  This proves (MA-MO.1506) for complex tests too.

\emph{2. The actual scalar functional yields a strong Sobolev boundary value.}
Use the fixed nonnegative compact $g_0,h_0,v_0$ of
(MA-MO.1204)--(MA-MO.1205).  Thus $g_0$ is supported
strictly inside $(1,2)$ and
\[
 k_0=C_{g_0^\vee}v_0\ge0,\qquad k_0\ne0,\qquad
 \operatorname{supp}k_0\subset J\subset(r_\phi-\log2,r_\phi),
\tag{MA-MO.1507}
\]
where $\phi$ has one nonzero sign on $J$ and
$r_\phi=\sup\operatorname{supp}\phi$.
For $T\in\mathbb R$ set
$g_T=e^{-iT\,\cdot}g_0$, $h_T=e^{-iT\,\cdot}h_0$ and
$v_T=e^{iT\,\cdot}v_0$.
Let $\ell_T$ be $\int v_Ty$ in the nontrivial case and
the actual functional $\Lambda_{v_T}$ of (MA-MO.1304)
in the rational case.  In either case
$\sup_T\|\ell_T\|<\infty$.  In the latter case this
includes the two compact polar-section integrals.

The bounded complex-linear functional $\ell_T T_{g_T}$ on
$L^2(w\,dx)$ has a bilinear representation
\[
       \ell_T(T_{g_T}y)=\int y(x)b_T(x)\,dx,\qquad
       b_T\in L^2(w^{-1}dx),\qquad
       \|b_T\|_{L^2(w^{-1})}\le C\|T_{g_T}\|.
\tag{MA-MO.1508}
\]
Complex conjugation is absorbed in $b_T$ by the Riesz
representation; the displayed pairing remains bilinear.
With $U_th(x)=h(x-t)$, the actual translated-source response is
\[
\begin{aligned}
 f_T(t)&=\ell_T(T_{g_T}^0U_th_T)
        =\int h_T(x-t)b_T(x)\,dx\\
       &=C_{h_T^\vee}b_T(t).
\end{aligned}
\tag{MA-MO.1509}
\]
The last equality has the indicated reflected convolution
kernel, without a conjugate.  For shifts in the fixed compact
support of $h_0^\vee$, $w(t-y)^{-1}/w(t)^{-1}$ is uniformly
bounded.  Minkowski's inequality therefore gives the weighted
Young estimate
\[
       \|f_T\|_{L^2(w^{-1})}
          \le C\|b_T\|_{L^2(w^{-1})}
          \le C'\|T_{g_T}\|.
\tag{MA-MO.1510}
\]
All constants are independent of $T$; the modulation changes
only the absolute phase of the convolution kernel.

Choose one $t_0>0$ satisfying the fixed support conditions
in (MA-MO.561).  Put
\[
       \eta_T(t)=\mathbf1_{[t_0,\infty)}(t)
                      e^{-(a+iT)t}f_T(t).
\tag{MA-MO.1511}
\]
Since
$w(t)^{-1}=e^{-2at}(1+e^{-4at})^{-1}(1+t^2)^{-r}$ for $t>0$,
\[
 \|\eta_T\|_{L^2((1+t^2)^{-r}dt)}
           \le\sqrt2\,\|f_T\|_{L^2(w^{-1})}
           \le C\|T_{g_T}\|.
\tag{MA-MO.1512}
\]
Consequently the actual half-line Laplace integral satisfies
\[
\begin{aligned}
 \mathfrak F_T(a+\epsilon+i(T+v))
   &=\int_{t_0}^{\infty}
             e^{-[a+\epsilon+i(T+v)]t}f_T(t)\,dt\\
   &=\sqrt{2\pi}\,\mathcal F_-
                      (e^{-\epsilon t}\eta_T)(v),\\
 \mathfrak F_T(a+\epsilon+i(T+\cdot))
   &\longrightarrow \sqrt{2\pi}\,\mathcal F_-\eta_T
                           \quad\text{strongly in }H^{-r}.
\end{aligned}
\tag{MA-MO.1513}
\]
Here $\mathcal F_-$ is the unitary minus Fourier transform.
The Sobolev norm on the second line is exactly the physical
$L^2((1+t^2)^{-r}dt)$ norm, up to $\sqrt{2\pi}$.
Dominated convergence proves the last line and the uniform
bound $C\|T_{g_T}\|$ for $0<\epsilon\le\epsilon_0$.
Absolute Laplace convergence and its holomorphy for
$\Re w>a$ follow either by weighted Cauchy--Schwarz or
the compact-source translation estimate.  This uses the
assumed bounded actual column, not a proposed inverse
multiplier on arbitrary data.

\emph{3. Exact noncancellation recovers the localized reciprocal.}
For $w=a+\epsilon+i(T+v)$ the original wide Euler calculation,
continued from $\Re w>1/2$ using (MA-MO.1501), is
\[
 \mathfrak F_T(w)
    =H_T(w)V_T(-w)
           -\frac{H_T(w)B_T(w)}{A_\eta(w)}-\mathcal R_T(w).
\tag{MA-MO.1514}
\]
The functions and actual compact remainder are exactly those
in (MA-MO.558)--(MA-MO.567).  In the rational case use
(MA-MO.1311)--(MA-MO.1313): first clear the zeta pole,
continue that identity, and remove the reciprocal's harmless
zero at $w=1/2$.  No output polar coordinate is discarded.
The Laplace function in (MA-MO.1514) is (MA-MO.1513), because
the source identity is valid on every compact translated
input and the wide Laplace integral is absolutely convergent.

Every nonunit ideal, or integer, shifts $J$ past $r_\phi$.
Thus only its unit term contributes to $B_T$.  Modulation
cancels exactly and gives
\[
\begin{aligned}
 H_T(w)&=H_0(a+\epsilon+iv)
             :=\int h_0(x)e^{(a+\epsilon+iv)x}\,dx,\\
 B_T(w)&=B_0(a+\epsilon+iv)
             :=\int_J\phi(x)k_0(x)e^{-(a+\epsilon+iv)x}\,dx,\\
 V_T(-w)&=V_0(-a-\epsilon-iv)
             :=\int_Jk_0(x)e^{-(a+\epsilon+iv)x}\,dx.
\end{aligned}
\tag{MA-MO.1515}
\]
At $v=0$ the first two have fixed nonzero signs and absolute
values uniformly bounded below for $0\le\epsilon\le\epsilon_0$.
Compactness therefore gives $v_*>0$ with
\[
 |U_\epsilon(v)|\ge c>0,\qquad
 U_\epsilon(v)=H_0(a+\epsilon+iv)B_0(a+\epsilon+iv),
 \qquad |v|\le v_* .
\tag{MA-MO.1516}
\]
All derivatives of $U_\epsilon^{-1}$ on a smaller fixed
interval are uniformly bounded and converge as $\epsilon\downarrow0$.
On the compact double-integration region defining $\mathcal R_T$,
the ideal sums are finite with uniformly bounded absolute
values.  Every fixed $v$ or $\epsilon$ derivative only inserts
a bounded power of a compact integration variable.  Hence
$\mathcal R_T(a+\epsilon+i(T+v))$ has uniform bounds in every
fixed smooth seminorm on this rectangle, and converges there
smoothly for each $T$ as $\epsilon\downarrow0$.

For any fixed $\chi_*\in C_c^\infty((-v_*,v_*))$, (MA-MO.1514)
now gives the actual identity
\[
\begin{aligned}
 \frac{\chi_*(v)}{A_\eta(a+\epsilon+i(T+v))}
 ={}&\frac{\chi_*(v)}{U_\epsilon(v)}
  \bigl[H_0(a+\epsilon+iv)V_0(-a-\epsilon-iv)\\
 &\quad-\mathcal R_T(a+\epsilon+i(T+v))
            -\mathfrak F_T(a+\epsilon+i(T+v))\bigr].
\end{aligned}
\tag{MA-MO.1517}
\]
Multiplication by a compact smooth function
is bounded on $H^{-r}$; its operator norm is controlled by
finitely many derivatives, by duality from (MA-MO.716).
The multipliers here converge in those norms.  Equations
(MA-MO.1513)--(MA-MO.1517) prove a strong $H^{-r}$ limit and
\[
 \sup_{0<\epsilon\le\epsilon_0}
 \left\|\frac{\chi_*(\,\cdot\,)}
 {A_\eta(a+\epsilon+i(T+\,\cdot\,))}\right\|_{H^{-r}}
       \le C(1+\|T_{g_T}\|)
       \le C'(2+|T|)^N .
\tag{MA-MO.1518}
\]
The final inequality is (MA-MO.1506), since
$\|g_T\|_{C^N}\le C_N(2+|T|)^N$.
Choose a uniform partition with sufficiently short supports
inside these windows and take $T=t_n$.  Frequency translation
is an isometry on $H^{-r}$, so this is (MA-MO.1502).

For the nontrivial sector, the positive tests have so far
given only the $\eta$ right boundary.  The actual source
reflection identity is
\[
 R\mathscr D_{\eta,g}^{\,\mathrm{raw}}R
       =\mathscr D_{\bar\eta,g^\vee}^{\,\mathrm{raw}},
                         \qquad Ry(x)=y(-x).
\tag{MA-MO.1519}
\]
It is (MA-MO.576), including the original Fourier/dual phases;
$R$ is an isometry of the sum norm.  The negative-support
part of (B) therefore supplies the positive-support dual
family, and the same argument supplies its right-boundary
condition.  No dual estimate is inferred from the positive
$\eta$ family alone.  The rational multiplier is self-dual,
so its original negative-column hypothesis already suffices.
This proves (B)$\Rightarrow$(L).

\emph{4. The boundary condition supplies the actual source inverse.}
First note that the partition convention is immaterial.
Given a second uniform smooth partition, multiply its $m$th
term by $\sum_n\chi_n=1$.  Only a bounded number of comparable
windows occur.  Their smooth multiplier norms are uniform,
and their indices have comparable height.  Thus both the
polynomial estimate and each strong limit in (MA-MO.1502)
pass to the second partition.  Refinement of either partition
if necessary gives the same argument for different fixed
window lengths.

Put $j_{\omega,+}=\sum_n d_{\omega,n,0}$.  This is a tempered
distribution: localized Schwartz tests decrease more rapidly
than the polynomial norms in (MA-MO.1502).  The same dominated
sum gives, for each $f\in\mathscr S$,
\[
\begin{gathered}
 A_\omega(a+\epsilon+i\,\cdot\,)^{-1}f
       \longrightarrow j_{\omega,+}f
                             \quad\text{strongly in }H^{-r},\\
 A_\omega(a+i\,\cdot\,)j_{\omega,+}=1 .
\end{gathered}
\tag{MA-MO.1520}
\]
The second equality follows from the identity at $\epsilon>0$
and the polynomial derivative bounds for the forward factor,
before taking the boundary limit.  It is not an arbitrary
principal-value prescription.  The first assertion also holds
when $f_\epsilon\to f$ in Schwartz space, and after any fixed
polynomial multiplication of $f$.

For $y\in\mathcal E$, let $F_{\omega,y}$ be the original
wide inverse.  At every fixed $b>a$ sufficiently close to $a$
the closed-strip premise supplies a positive margin, so the
already proved wide-source shift identifies
$V_bF_{\omega,y}=A_\omega(b+i\,\cdot\,)^{-1}V_by$.
For this fixed smooth input, $V_by\to V_ay$ in Schwartz space.
Equation (MA-MO.1520) therefore identifies the boundary of
that same original function, and all its fixed derivatives:
\[
 V_aF_{\omega,y}=j_{\omega,+}V_ay,\qquad
 D_x^j(e^{ax}F_{\omega,y})
                   \in L^2((1+x^2)^{-r}dx)\quad(j\ge0).
\tag{MA-MO.1521}
\]
The boundary convergence can also be read distributionally
in $x$: the original function is fixed on every compact
interval, so the limit cannot choose a different inverse.
Local Sobolev bounds give a polynomial bound for
$e^{ax}F_{\omega,y}$ and its derivatives on the left.
On the right it has the original rapid weighted bounds
at every wider line.  Thus $qF_{\omega,y}$ belongs to the
original local source, including the structural finite
vectors.  In the rational case it is precisely the actual
even additive Schwartz vector (MA-MO.714).

\emph{5. Dominated nuclear assembly retains the original matrix.}
We recall the needed estimate rather than infer boundedness
from a scalar boundary value.  If $d\in H^{-r}$ has compact
support and $L\in\mathscr S(\mathbb R^2)$, then
\[
 T_{d,L}f(t)=\langle d(u),L(t,u)f(u)\rangle
       \in\mathcal S_1(H^r,H^R)\quad(R\ge0).
\tag{MA-MO.1522}
\]
Expand $L$ in output Hermite functions.  Their $H^R$ norms
grow polynomially, while all multiplier seminorms of their
input coefficients decrease faster than any power.  The
$H^{-r}$--$H^r$ pairing and (MA-MO.716) give an absolutely
summable rank-one series.  This is the proof of (MA-MO.715)
and shows continuity in $d$ and finitely many Schwartz
seminorms of $L$.  Applied to the $n$th localized distribution,
Schwartz decrease in its support gives, for every $M$,
\[
 \|T_{d_{\omega,n,\epsilon},L}\|_{\mathcal S_1(H^r,H^R)}
          \le C_{L,M}(2+|n|)^{D-M}.
\tag{MA-MO.1523}
\]
Choose $M>D+2$.  Each summand converges in trace norm by
the strong limit in (MA-MO.1502); dominated summation proves
convergence of the entire series.  This remains true when
$L_\epsilon\to L_0$ in Schwartz space.  No uniform local
boundary-convergence modulus is used.

For completeness the two nonsingular source kernels are,
with $z=b+it,w=b+iu$, $G=\mathcal Mg$ and
$P(\xi)=\int e^{\xi x}\phi'(x)\,dx$,
\[
\begin{aligned}
 L_{\omega,b,g}(t,u)
   &=\frac{G(z)P(z-w)}{2\pi}
                         \frac{A_\omega(z)-A_\omega(w)}{z-w},\\
 L^{\mathrm{cross}}_{\eta,b,g}(t,u)
   &=\frac{G(z)A_\eta(z)m_\eta(z)}{2\pi}
                                      \frac{P(-z-w)}{-z-w}.
\end{aligned}
\tag{MA-MO.1524}
\]
These are exactly (MA-MO.718)--(MA-MO.719).  The first
is paired with the $\omega$ reciprocal and the crossed
one with the dual reciprocal.  The divided difference
is removed before distributional pairing; in the second
line $\Re(-z-w)\le-2a$.  Forward polynomial derivative
bounds, rapid decrease of $G$ and $P$, and the actual
local multiplier $m_\eta$ with its conductor, different
and finite phases show Schwartz convergence as $b\downarrow a$.
No reciprocal estimate is hidden in these nonsingular factors.

Apply (MA-MO.1523) and the exact identification
$H_{a,s}\simeq H^r$ of (MA-MO.717).  With the resulting
operators $K_{\eta,a}^\partial(g),O_{\eta,a}^\partial(g)$,
the first actual row is
\[
 \mathscr R_{\eta,a}^\partial(g)y
       =K_{\eta,a}^\partial(g)y-C_gM_\phi y
                              +O_{\eta,a}^\partial(g)Ry .
\tag{MA-MO.1525}
\]
The middle map is trace class from the sum space to
$H_{a,s}$ by its physical half-line kernel (MA-MO.332).
The other two are trace class by (MA-MO.1523).
For $y\in\mathcal E$, pass the exact original $b>a$ source
identity to the boundary using (MA-MO.1521), the strong
input convergence, and trace-norm kernel convergence.
The original smooth theta output is fixed during this
limit.  Thus (MA-MO.1525) agrees with the original row,
without an assumed bounded embedding between whole
varying-width spaces.

In the nontrivial sector the reflected dual row agrees
on $\mathcal E$ by actual source Poisson summation.
The paired map into $H_{a,s}\oplus H_{-a,s}$ is trace
class.  Its range lies in the closed diagonal copy of
$H$, first on the dense original core and hence everywhere.
The inverse diagonal isometry gives the required actual
trace-class endomorphism on the inherited sum norm.

\emph{6. The two rational coordinates and the remaining implications.}
In the rational case put $p_a(x)=\phi(x)e^{(1/2-a)x}$.
It is Schwartz.  Both actual cutoff moments extend boundedly:
\[
\begin{aligned}
 \lambda_0(y)&=\sum_n\langle d_{n,0}(u),
                         \widehat p_a(-u)(V_ay)(u)\rangle,\\
 \lambda_1(y)&=\lambda_0(Ry).
\end{aligned}
\tag{MA-MO.1526}
\]
The localized multiplier norms of $\widehat p_a(-u)$ decrease
rapidly, so (MA-MO.1502) makes this a norm-convergent sum
of bounded functionals.  Bilinear unitary Parseval on
$b>a$, followed by (MA-MO.1520), identifies them with the
original physical moments, with no extra Fourier constant.
The left bounds in (MA-MO.1521), $a<1/2$, and the original
right bounds make $\mathcal MF_y(1/2)$ absolutely convergent.
The original zeta inverse has a zero there, so it is zero.
Accordingly the actual source has
$h_y(0)=-\lambda_1(y)$ and $\int h_y=-\lambda_0(y)$.

Retain the fixed actual sections and $e_0,e_1$ of
(MA-MO.725).  The original diagonal and polar coordinates are
\[
\begin{aligned}
 T_g^\partial y={}&K_a^\partial(g)y-C_gM_\phi y
                       +O_a^\partial(g)Ry\\
 &+G(-1/2)\lambda_1(y)e_0+G(1/2)\lambda_0(y)e_1,\\
 \nu_g(y)={}&(-G(-1/2)\lambda_1(y),-G(1/2)\lambda_0(y)),\\
 \mathscr D_g(y,c)={}&(T_g^\partial y,\nu_g(y))
                                  +\sum_{j=0}^1c_jC_gDs_j .
\end{aligned}
\tag{MA-MO.1527}
\]
These are precisely the original polar signs and finite
columns.  Every term is trace class or finite rank; the
compatible diagonal argument is unchanged after the
displayed polar subtraction.  This proves (S) on $\mathcal W$.

All estimates in (MA-MO.1523)--(MA-MO.1527) involve only
finitely many derivatives of $g$ on a fixed support.
Thus the claimed trace-norm continuity holds.  The original
smooth cores are dense, proving uniqueness.  We have proved
(L)$\Rightarrow$(S), and (S)$\Rightarrow$(B) is immediate,
including restriction to the isometric rational negative
column.  Real and complex versions are equivalent by
Step 1.

Finally choose a nonzero compact test.  The necessary
actual boundary obstructions MA-MO-025 and MA-MO-026
applied to the just-constructed full matrix force
$s>2m-1$ for every boundary zero.  Therefore full individual
retention is a conclusion, not an unproved input to the
sufficiency argument.  MA-MO-028 supplies the actual ideal
trace; with full retention, the relative-trace and
determinant argument of MA-MO-029 applies verbatim.
\end{proof}

\paragraph{Meaning and limits of the criterion.}
The condition (L) contains no cutoff $\phi$.  Therefore universal
boundedness for one normalized smooth step implies the full
trace-class conclusion for every such step, with constants
allowed to depend on the step.  This follows by applying the
proved necessity for the first step and sufficiency for the
second; it asserts no cutoff independence for one isolated test.

This theorem removes the gap between the particular sufficient
cluster budgets and a necessary analytic boundary condition.
MA-MO-030 and MA-MO-035 prove (MA-MO.1502) for their stated
geometries.  The present equivalence does not show that full
individual jet retention alone implies it.  Its local
$H^{-s/2}$ norm permits oscillation and ordinate separation;
it is not replaced by the stronger nonoscillatory minimum-depth
budget.  Neither a hypothetical bad cluster nor failure of
that sufficient budget is asserted to occur.

The conclusion is for a universal fixed-support test family.
It makes no equivalence claim for one predetermined pair
$(g,\phi)$, for an unrestricted one-sided inverse composition,
or for the distinct whole-quotient return cost.  In particular
it does not decide RF002/RS003 or the PR cost gate merely by
rewriting them.  Under A the rational and Dirichlet fixed-gap
branch supplies the criterion; M supplies the indicated
Hecke branch.  Those zero-location inputs remain explicit.
No independent arithmetic deficiency vanishing in FT002
has been proved here.

\paragraph{Source and framework placement.}
The genuine source definitions, dual reflection and wide Euler
identity are MA-MO-024/026, with primary Meyer (2005) locations
in the supplied \texttt{main.tex}, lines 2630--2675,
3592--3707 and 3808--4043.  The exact polar completion is
MA-MO-016; the bilinear functional and universal-test
interfaces are MA-MO-033/034.  The distributional kernel
lemma, source comparison and compatible rows are
MA-MO-027/030/035.  The new necessity is the actual
weighted Young and strong boundary argument
(MA-MO.1508)--(MA-MO.1518), not a scalar trace assignment.
The unconditional fixed-strip forward estimates are
MT-M18.5--9.

In the inherited cote 1 placement, BR-TR-01/02,
scans 079--080 and 111--112, the finite-length
$\mathbb C[z]$-modules encode local divisor multiplicities
and their exterior-power determinants.  The strong
boundary topology and the actual nuclear series above
are additional analytic structure.  Finite-length
$K_0$ classes alone do not supply their norms, their
arithmetic validity, or a new positivity statement.
